% ===================================================================
% 《测圆海镜》卷四至卷六 —— 文言文↔白话文 对照本
% Ceyuan Haijing Volumes 4-6: Classical↔Modern Chinese Bilingual Edition
% ===================================================================
\documentclass[10pt,a4paper]{article}

%% -------- Core Packages --------
\usepackage{fontspec}
\usepackage{polyglossia}
\usepackage{amsmath,amssymb,amsthm}
\usepackage{geometry}
\usepackage{graphicx}
\usepackage{xcolor}
\usepackage{titlesec}
\usepackage{fancyhdr}
\usepackage{enumitem}
\usepackage{array,booktabs,longtable,multirow,colortbl}
\usepackage{hyperref}
\usepackage{microtype}
\usepackage{tocloft}
\usepackage{xeCJK}
\usepackage{paracol}

%% -------- Page Geometry (wider for two-column parallel text) --------
\geometry{
  paperwidth=210mm,paperheight=297mm,
  left=15mm,right=15mm,top=25mm,bottom=25mm,
  headheight=14pt,footskip=30pt
}

%% -------- Column Separation for paracol --------
\columnsep=8mm

%% -------- Language Setup --------
\setmainlanguage{english}

%% -------- Font Configuration --------
\setmainfont{Noto Serif}[
  UprightFont=* Regular,
  BoldFont=* Bold,
  ItalicFont=* Italic,
  BoldItalicFont=* Bold Italic
]
\setsansfont{Noto Sans}[
  UprightFont=* Regular,
  BoldFont=* Bold,
  ItalicFont=* Italic
]
\setmonofont{DejaVu Sans Mono}[Scale=MatchLowercase]

%% CJK Support via xeCJK — explicit path as required
\setCJKmainfont[Path=fonts/noto-cjk/]{NotoSansCJKsc-Regular.otf}
\setCJKsansfont[Path=fonts/noto-cjk/]{NotoSansCJKsc-Regular.otf}
\setCJKmonofont[Path=fonts/noto-cjk/]{NotoSansCJKsc-Regular.otf}

%% -------- Colors --------
\definecolor{titlecolor}{RGB}{45,55,72}
\definecolor{sectioncolor}{RGB}{55,65,81}
\definecolor{notecolor}{RGB}{120,53,15}
\definecolor{rulecolor}{RGB}{180,160,140}
\definecolor{volcolor}{RGB}{80,40,20}
\definecolor{leftbg}{RGB}{252,250,245}
\definecolor{rightbg}{RGB}{245,250,252}
\definecolor{leftborder}{RGB}{160,140,100}
\definecolor{rightborder}{RGB}{100,140,160}

%% -------- Section Formatting --------
\titleformat{\section}
  {\normalfont\Large\bfseries\color{volcolor}}
  {\thesection}{1em}{}
\titleformat{\subsection}
  {\normalfont\large\bfseries\color{sectioncolor}}
  {\thesubsection}{1em}{}
\titleformat{\subsubsection}
  {\normalfont\normalsize\bfseries\color{sectioncolor}}
  {\thesubsubsection}{1em}{}

%% -------- Header/Footer --------
\pagestyle{fancy}
\fancyhf{}
\fancyhead[L]{\small\textsc{Ceyuan Haijing — Classical↔Modern Chinese}}
\fancyhead[R]{\small\thepage}
\renewcommand{\headrulewidth}{0.4pt}
\renewcommand{\footrulewidth}{0pt}

%% -------- Custom Commands --------
\newcommand{\longline}{\noindent\rule{\textwidth}{0.4pt}}
\newcommand{\shortline}{\noindent\rule{0.5\textwidth}{0.4pt}\par}

%% Problem number command with Chinese numerals
\newcounter{probnum}
\newcommand{\chnum}[1]{%
  \ifcase#1 零\or 一\or 二\or 三\or 四\or 五\or 六\or 七\or 八\or 九\or 十\or
    十一\or 十二\or 十三\or 十四\or 十五\or 十六\or 十七\or 十八\fi}

%% Bilingual problem environment
%% LEFT = Classical Chinese, RIGHT = Modern Chinese
\newcommand{\biproblem}[2]{%
  \par\noindent
  \begin{minipage}[t]{0.48\textwidth}
    \textbf{第\chnum{\value{probnum}}问：}#1
  \end{minipage}%
  \hfill\vrule\hfill
  \begin{minipage}[t]{0.48\textwidth}
    \textbf{第\chnum{\value{probnum}}题：}#2
  \end{minipage}
  \par\vspace{0.3cm}
}

\newcommand{\bian answer}[2]{%
  \par\noindent
  \begin{minipage}[t]{0.48\textwidth}
    \textbf{答曰：}#1
  \end{minipage}%
  \hfill\vrule\hfill
  \begin{minipage}[t]{0.48\textwidth}
    \textbf{答：}#2
  \end{minipage}
  \par\vspace{0.2cm}
}

\newcommand{\biworking}[2]{%
  \par\noindent
  \begin{minipage}[t]{0.48\textwidth}
    \textbf{草曰：}#1
  \end{minipage}%
  \hfill\vrule\hfill
  \begin{minipage}[t]{0.48\textwidth}
    \textbf{演算：}#2
  \end{minipage}
  \par\vspace{0.2cm}
}

\newcommand{\bimethod}[2]{%
  \par\noindent
  \begin{minipage}[t]{0.48\textwidth}
    \textbf{法曰：}#1
  \end{minipage}%
  \hfill\vrule\hfill
  \begin{minipage}[t]{0.48\textwidth}
    \textbf{算法：}#2
  \end{minipage}
  \par\vspace{0.3cm}
}

%% Simplified combined command for problems with integrated method
\newcommand{\problemmethod}[2]{%
  \par\noindent
  \begin{minipage}[t]{0.48\textwidth}
    #1
  \end{minipage}%
  \hfill\vrule\hfill
  \begin{minipage}[t]{0.48\textwidth}
    #2
  \end{minipage}
  \par\vspace{0.3cm}
}

%% Section headers bilingual
\newcommand{\bisectionheader}[4]{%
  \begin{center}
  {\Large\bfseries\color{volcolor} #1}\\[0.3cm]
  {\normalsize\color{sectioncolor} #2}
  \end{center}
  \vspace{0.3cm}
  \noindent
  \begin{minipage}[t]{0.48\textwidth}
    \textbf{【题意总说】}#3
  \end{minipage}%
  \hfill\vrule\hfill
  \begin{minipage}[t]{0.48\textwidth}
    \textbf{【题意概述】}#4
  \end{minipage}
  \par\vspace{0.3cm}
}

%% -------- Hyperref Setup --------
\hypersetup{
  colorlinks=true,
  linkcolor=sectioncolor,
  citecolor=sectioncolor,
  urlcolor=sectioncolor,
  pdfencoding=auto,
  pdftitle={测圆海镜卷四至卷六 文言文↔白话文对照本},
  pdfauthor={李冶 (Li Ye)}
}

%% -------- TOC Depth --------
\setcounter{tocdepth}{3}
\setcounter{secnumdepth}{3}

%% ===================================================================

\begin{document}

%% ===================================================================
%% TITLE PAGE
%% ===================================================================
\begin{titlepage}
\begin{center}
\vspace*{2cm}

{\fontsize{14}{18}\selectfont 钦定四库全书}

\vspace{1.5cm}

{\fontsize{12}{16}\selectfont 子部}

\vspace{2cm}

\rule{0.8\textwidth}{1.5pt}

\vspace{1cm}

{\fontsize{26}{34}\selectfont\bfseries\color{titlecolor} 测圆海镜}

\vspace{0.3cm}

{\fontsize{16}{22}\selectfont 卷四至卷六}

\vspace{0.3cm}

{\fontsize{13}{17}\selectfont\color{notecolor} 文言文 \raisebox{-0.5ex}{$\leftrightarrow$} 白话文 对照本}

\vspace{1cm}

\rule{0.8\textwidth}{1.5pt}

\vspace{2cm}

{\fontsize{14}{18}\selectfont 元\quad 李冶\quad 撰}

\vspace{0.5cm}

{\fontsize{11}{15}\selectfont\color{sectioncolor} 白话文译解}

\vspace{3cm}

\begin{tabular}{cc}
\textbf{左栏} & \textbf{右栏} \\
\textit{文言文（原文）} & \textit{白话文（今译）} \\
\end{tabular}

\vfill

{\small 依《钦定四库全书》本 \quad 现代汉语译注}

\end{center}
\end{titlepage}

%% ===================================================================
%% TABLE OF CONTENTS
%% ===================================================================
\tableofcontents
\newpage

%% ===================================================================
%% INTRODUCTION (BILINGUAL)
%% ===================================================================
\section*{引言}
\addcontentsline{toc}{section}{引言}

\noindent
\begin{minipage}[t]{0.48\textwidth}
\textbf{【原文】}

《测圆海镜》十二卷，元李冶撰。冶字仁卿，号敬斋，真定栾城人。
金末进士，入元官翰林学士。此书乃其所著天元术之杰作，以勾股测圆
之法，设问答以明天元一之术。全书凡一百七十问，皆设城池圆形，以
人物行步之数推求圆径。

本编收录卷四、卷五、卷六三卷：
\begin{itemize}[leftmargin=1em]
  \item \textbf{卷四}：底勾一十七问
  \item \textbf{卷五}：大股一十八问  
  \item \textbf{卷六}：大勾一十八问
\end{itemize}

书中设问之例：假令有圆城一座，甲乙二人各从城门出行，或东或南，
行若干步，遥望见之，或斜行相会。只知若干步数，余皆不知，以求
圆城之径。每问皆有\textbf{问}、\textbf{答}、\textbf{草}、\textbf{法}四部分。

\textbf{术语解释：}
\begin{itemize}[leftmargin=1em]
  \item \textbf{天元一}：设未知数为一，即现代代数之$x$
  \item \textbf{寄左}：将某式暂存一边，相当于方程之一边
  \item \textbf{相消}：两式相减，相当于移项合并
  \item \textbf{带分母}：某式含有分母未除
  \item \textbf{幂}：平方，即现代记号$x^2$
  \item \textbf{开方}：开平方，即求平方根
  \item \textbf{开立方}：开立方，即求立方根
  \item \textbf{实}：常数项
  \item \textbf{从}：一次项系数
  \item \textbf{廉}：二次项系数
  \item \textbf{隅}：三次项系数（立方方程中）
\end{itemize}
\end{minipage}%
\hfill\vrule\hfill
\begin{minipage}[t]{0.48\textwidth}
\textbf{【今译】}

《测圆海镜》共十二卷，由元代数学家李冶撰写。李冶字仁卿，号敬斋，
真定府栾城县人。金朝末年考中进士，元朝时任翰林学士。这部书是他
研究天元术的杰作，运用勾股定理测量圆形的方法，通过设置问题和
解答来阐明天元术的奥秘。全书共170道题，都假设有一座圆形城池，
根据人物行走的步数来推算城池的直径。

本版本收录第四、第五、第六卷：
\begin{itemize}[leftmargin=1em]
  \item \textbf{卷四}：底勾17题
  \item \textbf{卷五}：大股18题
  \item \textbf{卷六}：大勾18题
\end{itemize}

书中设题的模式：假设有一座圆形城池，甲、乙两人分别从城门出发，
或向东或向南行走若干步，远远望见对方，或斜向行走相会。只给定
某些步数，其余未知，要求求出圆形城池的直径。每道题都有
\textbf{问题}、\textbf{答案}、\textbf{演算}、\textbf{算法}四部分。

\textbf{术语解释：}
\begin{itemize}[leftmargin=1em]
  \item \textbf{天元一}：设未知数为1，相当于现代代数的$x$
  \item \textbf{寄左}：将某个算式暂存一边，相当于方程的一边
  \item \textbf{相消}：两个算式相减，相当于移项合并同类项
  \item \textbf{带分母}：某个算式含有分母尚未除尽
  \item \textbf{幂}：平方，即现代记号$x^2$
  \item \textbf{开方}：开平方，即求平方根
  \item \textbf{开立方}：开立方，即求立方根
  \item \textbf{实}：常数项
  \item \textbf{从}：一次项系数
  \item \textbf{廉}：二次项系数
  \item \textbf{隅}：三次项系数（立方方程中）
\end{itemize}
\end{minipage}

\vspace{0.5cm}
\longline

\newpage

%% ===================================================================
%% FORMAT NOTES
%% ===================================================================
\noindent
\begin{minipage}[t]{0.48\textwidth}
\textbf{【体例说明】}

左栏为《测圆海镜》原文，依四库全书本。
右栏为白话文翻译，力求准确通顺。

数学算式保持原书格式，以天元术排列
多项式。必要时加注说明。
\end{minipage}%
\hfill\vrule\hfill
\begin{minipage}[t]{0.48\textwidth}
\textbf{【格式说明】}

左栏为原文，依据《四库全书》版本。
右栏为白话文翻译，力求准确、通顺。

数学算式保持原书格式，采用天元术的
方式排列多项式。必要时添加注释说明。
\end{minipage}

\vspace{0.5cm}
\longline

\newpage

%% ===================================================================
%% VOLUME 4
%% ===================================================================
\section{测圆海镜卷四 —— 底勾一十七问}
\label{sec:vol4}

\bisectionheader
  {底勾一十七问}
  {十七道关于底勾的数学题}
  {卷四论「底勾」之术。所谓底勾，乃勾股形之最下勾股，与圆径相切之基本勾股形也。凡一十七问，皆以底勾为基，求圆城之径。}
  {卷四讨论「底勾」的算法。所谓底勾，是指勾股三角形中最下面的勾股边，是与圆直径相切的基本勾股三角形。共17道题，都以底勾为基础，求圆形城池的直径。}

\vspace{0.3cm}

%% Reset problem counter for Volume 4
\setcounter{probnum}{0}

%% --- Problem 1 ---
\stepcounter{probnum}
\biproblem
  {或问乙出南门东行不知步数而立，甲出北门东行二百步见之，就乙斜行二百七十二步与乙相会。问答同前。}
  {假设乙从南门出发向东行走若干步停下，甲从北门出发向东行走200步后看见乙，接着向乙斜行272步与乙相会。问题与答案同前。}

\bian answer
  {城径二百四十步。}
  {城池直径为240步。}

\biworking
  {识别得二行相减余即乙出南门东行数也。以甲东行减于就乙斜行余七十二步，以乘甲东行步，得一万四千四百步，又四之，得五万七千六百步为实，以平方开之，得二百四十步，即城径也。}
  {通过分析可知，两次行走的差就是乙出南门向东行走的步数。用甲向东行走的步数从斜行步数中减去，余72步，乘以甲向东行走的步数，得到14400步，再乘以4，得到57600步作为被开方数（实），开平方得到240步，即城池的直径。}

\bimethod
  {二行差数乘甲东行又四之，为平方实，得全径。}
  {两次行走的差乘以甲向东行走的步数再乘以4，作为平方开方的被开方数，求得全城直径。}

%% --- Problem 2 ---
\stepcounter{probnum}
\biproblem
  {或问乙出南门东行不知步数而立，甲出北门东行二百步见之。问答同前。}
  {假设乙从南门出发向东行走若干步停下，甲从北门出发向东行走200步后看见乙。问题与答案同前。}

\bian answer
  {城径二百四十步。}
  {城池直径为240步。}

\biworking
  {立天元一为城径，以减于二之甲东行步，得\raisebox{-0.5ex}{\scriptsize$\begin{array}{c}480\\\hline 1\end{array}$}为两个小差，以乙南行步乘之，得\raisebox{-0.5ex}{\scriptsize$\begin{array}{c}48000\\\hline 1\end{array}$}为城径幂，寄左。然后以天元幂\raisebox{-0.5ex}{\scriptsize$\begin{array}{c}0\\1\end{array}$}与左相消，得\raisebox{-0.5ex}{\scriptsize$\begin{array}{c|c|c}0&0&48000\\\hline 1&0&\end{array}$}以平方开之，得二百四十步，即城径也。}
  {设天元一（未知数）为城池直径，用两倍的甲向东行走步数减去它，得到$\begin{pmatrix}480\\1\end{pmatrix}$为两个小差，乘以乙向南行走的步数，得到$\begin{pmatrix}48000\\1\end{pmatrix}$作为城径的平方，暂存左边。然后将天元的平方$\begin{pmatrix}0\\1\end{pmatrix}$与左边相消，得到$\begin{pmatrix}0&0&48000\\1&0\end{pmatrix}$，开平方得到240步，即城池直径。}

\bimethod
  {半之乙南行步乘甲东行为实，半乙南行为从，一步常法，得半径。}
  {取乙南行步数的一半乘以甲向东行走的步数作为被开方数，取乙南行步数的一半作为一次项系数，二次项系数为1，开方求得半径。}

%% --- Problem 3 ---
\stepcounter{probnum}
\biproblem
  {或问乙从坤隅南行三百六十步，甲出北门东行二百步见之。问答同前。}
  {假设乙从坤隅（西南角）向南行走360步，甲从北门出发向东行走200步后看见乙。问题与答案同前。}

\bian answer
  {城径二百四十步。}
  {城池直径为240步。}

\biworking
  {立天元一为半径，减于乙东行，得\raisebox{-0.5ex}{\scriptsize$\begin{array}{c}120\\1\end{array}$}为小差。半乙南行步，得一百八十步，以乘小差，得\raisebox{-0.5ex}{\scriptsize$\begin{array}{c}21600\\\hline 1\end{array}$}为半径幂，寄左。然后以天元幂\raisebox{-0.5ex}{\scriptsize$\begin{array}{c}0\\1\end{array}$}与左相消，得\raisebox{-0.5ex}{\scriptsize$\begin{array}{c|c|c}0&0&21600\\\hline 1&0&\end{array}$}以平方开之，得一百二十步，倍之得全径。}
  {设天元一为半径，从乙向东行走的距离中减去，得到$\begin{pmatrix}120\\1\end{pmatrix}$为小差。取乙向南行走步数的一半，得180步，乘以小差，得到$\begin{pmatrix}21600\\1\end{pmatrix}$作为半径的平方，暂存左边。然后将天元的平方$\begin{pmatrix}0\\1\end{pmatrix}$与左边相消，得到$\begin{pmatrix}0&0&21600\\1&0\end{pmatrix}$，开平方得到120步，乘以2得到全城直径。}

\bimethod
  {两行步相乘为实，甲东行为从，乙为隅，得半径。}
  {两行的步数相乘作为被开方数，甲向东行走的步数为一次项系数，乙的行步数为二次项系数，开方求得半径。}

%% --- Problem 4 ---
\stepcounter{probnum}
\biproblem
  {或问乙从坤隅南行三百六十步，甲出北门东行二百步见之。问答同前。}
  {假设乙从坤隅（西南角）向南行走360步，甲从北门出发向东行走200步后看见乙。问题与答案同前。}

\bian answer
  {城径二百四十步。}
  {城池直径为240步。}

\biworking
  {立天元一为城径，以减于二之甲东行，余即乙出南门东行数也。}
  {设天元一为城池直径，用两倍的甲向东行走步数减去它，余数就是乙出南门向东行走的步数。}

\bimethod
  {以乙南行步乘甲东行幂，又四之为实，从空，乙行为廉，一步常法，得城径。}
  {用乙向南行走的步数乘以甲向东行走步数的平方，再乘以4作为被开方数，一次项为空，乙的行步数为二次项系数，三次项系数为1，开方求得城径。}

%% --- Problem 5 ---
\stepcounter{probnum}
\biproblem
  {或问乙出南门直行一百三十五步，甲出北门东行二百步见之。问答同前。}
  {假设乙从南门出发向南直走135步，甲从北门出发向东行走200步后看见乙。问题与答案同前。}

\bian answer
  {城径二百四十步。}
  {城池直径为240步。}

\biworking
  {立天元一为城径，加乙南行，得\raisebox{-0.5ex}{\scriptsize$\begin{array}{c}135\\1\end{array}$}为股率。其甲东行即勾率也。其乙南行为小股，以勾率乘之，合以大弦率除，今不受除，便以此为小勾，为母，寄股率。乃先以小弦乘大和，得下式，寄左。又以倍天元减斜步，得为小和，以乘大弦，得同数，与左相消。}
  {设天元一为城池直径，加上乙向南行走的步数，得到$\begin{pmatrix}135\\1\end{pmatrix}$为股率。甲向东行走的步数就是勾率。乙向南行走的步数为小股，用勾率乘以它，应当用大弦率来除，现在先不除，就把这作为小勾，作为分母，暂寄股率。先用小弦乘以大和，得到下式，暂存左边。又用两倍的天元减去斜行步数，得到小和，乘以大弦，得到相同的数，与左边相消。}

\bimethod
  {以乙南行步乘甲东行幂，又四之为实，从空，乙行为廉，一步常法，得城径。}
  {用乙向南行走的步数乘以甲向东行走步数的平方，再乘以4作为被开方数，一次项为空，乙的行步数为二次项系数，三次项系数为1，开方求得城径。}

\vspace{0.3cm}
\noindent
\begin{minipage}[t]{0.48\textwidth}
\textit{（以下各问草法类推，皆立天元一为半径或城径，以各条件建立方程，与左式相消，开方得数。）}
\end{minipage}%
\hfill\vrule\hfill
\begin{minipage}[t]{0.48\textwidth}
\textit{（以下各题演算方法类似，都设天元一为半径或城径，根据各题条件建立方程，与左边算式相消，开方得到答案。）}
\end{minipage}

\vspace{0.3cm}

%% --- Problem 6 ---
\stepcounter{probnum}
\biproblem
  {或问乙从坤隅南行三百六十步，甲出北门东行二百步见之，就乙斜行二百七十二步与乙相会。问答同前。}
  {假设乙从坤隅（西南角）向南行走360步，甲从北门出发向东行走200步后看见乙，接着向乙斜行272步与乙相会。问题与答案同前。}

\bimethod
  {两行步相乘倍之为实，乙南行为从，一步常法，得半径。}
  {两行步数相乘再乘以2作为被开方数，乙向南行走的步数为一次项系数，二次项系数为1，开方求得半径。}

%% --- Problem 7 ---
\stepcounter{probnum}
\biproblem
  {或问乙出南门直行一百三十五步，甲出北门东行二百步见之。问答同前。}
  {假设乙从南门出发向南直走135步，甲从北门出发向东行走200步后看见乙。问题与答案同前。}

\bimethod
  {以乙南行乘甲东行幂为实，从空，乙行为廉，一步常法，得城径。}
  {用乙向南行走的步数乘以甲向东行走步数的平方作为被开方数，一次项为空，乙的行步数为二次项系数，三次项系数为1，开方求得城径。}

%% --- Problem 8 ---
\stepcounter{probnum}
\biproblem
  {或问乙出南门直行一百三十五步，甲出北门东行二百步见之。问答同前。}
  {假设乙从南门出发向南直走135步，甲从北门出发向东行走200步后看见乙。问题与答案同前。}

\bimethod
  {倍乙南行乘甲东行幂为实，从空，倍乙行为廉，一步常法，得城径。}
  {乙向南行走步数的2倍乘以甲向东行走步数的平方作为被开方数，一次项为空，乙行步数的2倍为二次项系数，三次项系数为1，开方求得城径。}

%% --- Problem 9 ---
\stepcounter{probnum}
\biproblem
  {或问乙出南门直行一百三十五步，甲出北门东行二百步，就乙斜行二百七十二步与乙相会。问答同前。}
  {假设乙从南门出发向南直走135步，甲从北门出发向东行走200步，接着向乙斜行272步与乙相会。问题与答案同前。}

\bimethod
  {倍乙南行乘甲东行幂为实，甲东行为从，倍乙行为廉，一步常法，得城径。}
  {乙向南行走步数的2倍乘以甲向东行走步数的平方作为被开方数，甲向东行走的步数为一次项系数，乙行步数的2倍为二次项系数，三次项系数为1，开方求得城径。}

%% --- Problem 10 ---
\stepcounter{probnum}
\biproblem
  {或问乙出南门直行一百三十五步，甲出北门东行二百步见之，就乙斜行二百七十二步与乙相会。问答同前。}
  {假设乙从南门出发向南直走135步，甲从北门出发向东行走200步后看见乙，接着向乙斜行272步与乙相会。问题与答案同前。}

\bimethod
  {两行相减余乘甲东行又四之为实，从空，两行差为廉，一步常法，得城径。}
  {两行步数相减的余数乘以甲向东行走步数再乘以4作为被开方数，一次项为空，两行步数的差为二次项系数，三次项系数为1，开方求得城径。}

%% --- Problem 11 ---
\stepcounter{probnum}
\biproblem
  {或问乙出南门直行一百三十五步，甲出北门东行二百步见之。问答同前。}
  {假设乙从南门出发向南直走135步，甲从北门出发向东行走200步后看见乙。问题与答案同前。}

\bimethod
  {以乙南行步乘甲东行幂，又四之为实，从空，乙行为廉，一步常法，得城径。}
  {用乙向南行走的步数乘以甲向东行走步数的平方，再乘以4作为被开方数，一次项为空，乙的行步数为二次项系数，三次项系数为1，开方求得城径。}

%% --- Problem 12 ---
\stepcounter{probnum}
\biproblem
  {或问乙出南门直行一百三十五步，甲出北门东行二百步见之。问答同前。}
  {假设乙从南门出发向南直走135步，甲从北门出发向东行走200步后看见乙。问题与答案同前。}

\bimethod
  {倍乙南行乘甲东行幂为实，从空，倍乙行为廉，一步常法，得城径。}
  {乙向南行走步数的2倍乘以甲向东行走步数的平方作为被开方数，一次项为空，乙行步数的2倍为二次项系数，三次项系数为1，开方求得城径。}

%% --- Problem 13 ---
\stepcounter{probnum}
\biproblem
  {或问乙出南门直行一百三十五步，甲出北门东行二百步，就乙斜行二百七十二步与乙相会。问答同前。}
  {假设乙从南门出发向南直走135步，甲从北门出发向东行走200步，接着向乙斜行272步与乙相会。问题与答案同前。}

\bimethod
  {两行相减余乘甲东行又四之为实，从空，两行差为廉，一步常法，得城径。}
  {两行步数相减的余数乘以甲向东行走步数再乘以4作为被开方数，一次项为空，两行步数的差为二次项系数，三次项系数为1，开方求得城径。}

%% --- Problem 14 ---
\stepcounter{probnum}
\biproblem
  {或问乙出南门直行一百三十五步，甲出北门东行二百步见之。问答同前。}
  {假设乙从南门出发向南直走135步，甲从北门出发向东行走200步后看见乙。问题与答案同前。}

\bimethod
  {以乙南行乘甲东行幂为实，从空，乙行为廉，一步常法，得城径。}
  {用乙向南行走的步数乘以甲向东行走步数的平方作为被开方数，一次项为空，乙的行步数为二次项系数，三次项系数为1，开方求得城径。}

%% --- Problem 15 ---
\stepcounter{probnum}
\biproblem
  {或问乙出南门直行一百三十五步，甲出北门东行二百步见之。问答同前。}
  {假设乙从南门出发向南直走135步，甲从北门出发向东行走200步后看见乙。问题与答案同前。}

\bimethod
  {倍乙南行乘甲东行幂为实，从空，倍乙行为廉，一步常法，得城径。}
  {乙向南行走步数的2倍乘以甲向东行走步数的平方作为被开方数，一次项为空，乙行步数的2倍为二次项系数，三次项系数为1，开方求得城径。}

%% --- Problem 16 ---
\stepcounter{probnum}
\biproblem
  {或问乙出南门直行一百三十五步，甲出北门东行二百步，就乙斜行二百七十二步与乙相会。问答同前。}
  {假设乙从南门出发向南直走135步，甲从北门出发向东行走200步，接着向乙斜行272步与乙相会。问题与答案同前。}

\bimethod
  {两行相减余乘甲东行又四之为实，从空，两行差为廉，一步常法，得城径。}
  {两行步数相减的余数乘以甲向东行走步数再乘以4作为被开方数，一次项为空，两行步数的差为二次项系数，三次项系数为1，开方求得城径。}

%% --- Problem 17 ---
\stepcounter{probnum}
\biproblem
  {或问乙出南门直行一百三十五步，甲出北门东行二百步见之。问答同前。}
  {假设乙从南门出发向南直走135步，甲从北门出发向东行走200步后看见乙。问题与答案同前。}

\bimethod
  {以乙南行步乘甲东行幂，又四之为实，从空，乙行为廉，一步常法，得城径。}
  {用乙向南行走的步数乘以甲向东行走步数的平方，再乘以4作为被开方数，一次项为空，乙的行步数为二次项系数，三次项系数为1，开方求得城径。}

\vspace{0.5cm}
\longline
\begin{center}
\textit{测圆海镜卷四终}
\end{center}

\newpage

%% ===================================================================
%% VOLUME 5
%% ===================================================================
\section{测圆海镜卷五 —— 大股一十八问}
\label{sec:vol5}

\bisectionheader
  {大股一十八问}
  {十八道关于大股的数学题}
  {卷五论「大股」之术。所谓大股，乃最大勾股形之股边，即通股也。以大股为主，配以诸勾股形之关系，求圆城之径。凡一十八问。}
  {卷五讨论「大股」的算法。所谓大股，是指最大勾股三角形的股边，即通股。以大股为主要对象，配合各种勾股三角形的关系，求圆形城池的直径。共18道题。}

\vspace{0.3cm}

%% Reset problem counter for Volume 5
\setcounter{probnum}{0}

%% --- Problem 1 ---
\stepcounter{probnum}
\biproblem
  {或问乙出南门直行一百三十五步，而立甲从乾隅南行六百步，斜望乙与城参相直。问答同前。}
  {假设乙从南门出发向南直走135步停下，甲从乾隅（西北角）向南行走600步，斜向遥望乙，视线与城池边缘恰好在一条直线上。问题与答案同前。}

\bian answer
  {城径二百四十步。}
  {城池直径为240步。}

\biworking
  {立天元一为半径，以二之减于甲南行，得\raisebox{-0.5ex}{\scriptsize$\begin{array}{c}600\\2\end{array}$}为大差也。以自之，得\raisebox{-0.5ex}{\scriptsize$\begin{array}{c}360000\\2400\\4\end{array}$}为大差幂也。置甲南行幂内减大差幂，而半之，得大弦也。内带大差为分母。又置甲南行幂内减大差幂，而半之，得大勾也。带大差分母。又以大差乘股六百步，併入和，得下式，为小和，以乘大弦。寄左。又以倍天元减斜步，得小和，以乘大弦，得同数，与左相消。开立方得一百二十步，即半径也。}
  {设天元一为半径，用甲向南行走步数减去两倍的天元，得到$\begin{pmatrix}600\\2\end{pmatrix}$为大差。将其平方，得到$\begin{pmatrix}360000\\2400\\4\end{pmatrix}$为大差的平方。将甲向南行走步数的平方减去大差的平方，再取一半，得到大弦。内含大差作为分母。又将甲向南行走步数的平方减去大差的平方，再取一半，得到大勾。也带大差分母。又用大差乘以股600步，合并入和，得到下式，作为小和，乘以大弦。暂存左边。又用两倍的天元减去斜行步数，得到小和，乘以大弦，得到相同的数，与左边相消。开立方得到120步，即半径。}

\bimethod
  {倍二行差内减甲南行步，复以乘甲南行步为实。倍二行差减甲南行步即甲南行也。四之甲南行步内减于甲南行余二百四十步，即城径也。}
  {两倍的两行差减去甲向南行走步数，再乘以甲向南行走步数作为被开方数。两倍的两行差减去甲向南行走步数就是甲向南行走步数。四倍的甲向南行走步数减去甲向南行走步数，余240步，即城池直径。}

%% --- Problem 2 ---
\stepcounter{probnum}
\biproblem
  {或问乙出南门直行一百三十五步，甲从乾隅南行六百步，斜望乙与城参相直。问答同前。}
  {假设乙从南门出发向南直走135步，甲从乾隅（西北角）向南行走600步，斜向遥望乙，视线与城池边缘恰好在一条直线上。问题与答案同前。}

\bian answer
  {城径二百四十步。}
  {城池直径为240步。}

\biworking
  {立天元一为半径，以减于甲南行步，得大差。置甲南行幂内减大差幂，而半之，得大弦也。又以大差乘股六百步，併入和，得下式，为小和，以乘大弦。寄左。又以倍天元减斜步，得小和，以乘大弦，得同数，与左相消。开立方得半径。}
  {设天元一为半径，从甲向南行走步数中减去，得到大差。将甲向南行走步数的平方减去大差的平方，再取一半，得到大弦。又用大差乘以股600步，合并入和，得到下式，作为小和，乘以大弦。暂存左边。又用两倍的天元减去斜行步数，得到小和，乘以大弦，得到相同的数，与左边相消。开立方得到半径。}

\bimethod
  {两行步相乘为实，从空，两行步为廉，一步常法，得城径。}
  {两行步数相乘作为被开方数，一次项为空，两行步数为二次项系数，三次项系数为1，开方求得城径。}

%% --- Problem 3 ---
\stepcounter{probnum}
\biproblem
  {或问乙出南门直行一百三十五步，甲从乾隅南行六百步望乙。问答同前。}
  {假设乙从南门出发向南直走135步，甲从乾隅（西北角）向南行走600步遥望乙。问题与答案同前。}

\bian answer
  {城径二百四十步。}
  {城池直径为240步。}

\biworking
  {立天元一为城径，以二之减于甲南行，余为大差。以自之得大差幂。置甲南行幂内减大差幂而半之，得大弦。内带大差为分母。又置甲南行幂内减大差幂而半之，得大勾。带大差分母。又以大差乘股六百步，併入和，得下式，为小和，以乘大弦。寄左。又以倍天元减斜步，得小和，以乘大弦，得同数，与左相消。开立方得一百二十步，即半径。}
  {设天元一为城池直径，用甲向南行走步数减去两倍的天元，余数为大差。将其平方得到大差的平方。将甲向南行走步数的平方减去大差的平方，再取一半，得到大弦。内含大差作为分母。又将甲向南行走步数的平方减去大差的平方，再取一半，得到大勾。也带大差分母。又用大差乘以股600步，合并入和，得到下式，作为小和，乘以大弦。暂存左边。又用两倍的天元减去斜行步数，得到小和，乘以大弦，得到相同的数，与左边相消。开立方得到120步，即半径。}

\bimethod
  {倍二行差内减甲南行步，复以乘甲南行步为实。倍二行差减甲南行步即甲南行也。四之甲南行步内减于甲南行余，即城径也。}
  {两倍的两行差减去甲向南行走步数，再乘以甲向南行走步数作为被开方数。两倍的两行差减去甲向南行走步数就是甲向南行走步数。四倍的甲向南行走步数减去甲向南行走步数，余数即城池直径。}

\vspace{0.3cm}
\noindent
\begin{minipage}[t]{0.48\textwidth}
\textit{（以下各问草法类推，皆依大股之术，以大差为主，建立方程。）}
\end{minipage}%
\hfill\vrule\hfill
\begin{minipage}[t]{0.48\textwidth}
\textit{（以下各题演算方法类似，都依照大股的算法，以大差为主要对象，建立方程。）}
\end{minipage}

\vspace{0.3cm}

%% --- Problem 4 ---
\stepcounter{probnum}
\biproblem
  {或问乙出南门直行一百三十五步，甲从乾隅南行六百步望乙。问答同前。}
  {假设乙从南门出发向南直走135步，甲从乾隅（西北角）向南行走600步遥望乙。问题与答案同前。}

\bimethod
  {两行步相乘为实，从空，两行步为廉，一步常法，得城径。}
  {两行步数相乘作为被开方数，一次项为空，两行步数为二次项系数，三次项系数为1，开方求得城径。}

%% --- Problem 5 ---
\stepcounter{probnum}
\biproblem
  {或问乙出南门直行一百三十五步，甲从乾隅南行六百步，斜望乙与城参相直。问答同前。}
  {假设乙从南门出发向南直走135步，甲从乾隅（西北角）向南行走600步，斜向遥望乙，视线与城池边缘恰好在一条直线上。问题与答案同前。}

\bimethod
  {倍二行差内减甲南行步，复以乘甲南行步为实，从空，倍二行差减甲南行步为廉，一步常法，得城径。}
  {两倍的两行差减去甲向南行走步数，再乘以甲向南行走步数作为被开方数，一次项为空，两倍的两行差减去甲向南行走步数作为二次项系数，三次项系数为1，开方求得城径。}

%% --- Problem 6 ---
\stepcounter{probnum}
\biproblem
  {或问乙出南门直行一百三十五步，甲从乾隅南行六百步望乙。问答同前。}
  {假设乙从南门出发向南直走135步，甲从乾隅（西北角）向南行走600步遥望乙。问题与答案同前。}

\bimethod
  {两行步相乘为实，从空，两行步为廉，一步常法，得城径。}
  {两行步数相乘作为被开方数，一次项为空，两行步数为二次项系数，三次项系数为1，开方求得城径。}

%% --- Problem 7 ---
\stepcounter{probnum}
\biproblem
  {或问乙出南门直行一百三十五步，甲从乾隅南行六百步，斜望乙与城参相直。问答同前。}
  {假设乙从南门出发向南直走135步，甲从乾隅（西北角）向南行走600步，斜向遥望乙，视线与城池边缘恰好在一条直线上。问题与答案同前。}

\bimethod
  {倍二行差内减甲南行步，复以乘甲南行步为实，从空，倍二行差减甲南行步为廉，一步常法，得城径。}
  {两倍的两行差减去甲向南行走步数，再乘以甲向南行走步数作为被开方数，一次项为空，两倍的两行差减去甲向南行走步数作为二次项系数，三次项系数为1，开方求得城径。}

%% --- Problem 8 ---
\stepcounter{probnum}
\biproblem
  {或问乙出南门直行一百三十五步，甲从乾隅南行六百步望乙。问答同前。}
  {假设乙从南门出发向南直走135步，甲从乾隅（西北角）向南行走600步遥望乙。问题与答案同前。}

\bimethod
  {两行步相乘为实，从空，两行步为廉，一步常法，得城径。}
  {两行步数相乘作为被开方数，一次项为空，两行步数为二次项系数，三次项系数为1，开方求得城径。}

%% --- Problem 9 ---
\stepcounter{probnum}
\biproblem
  {或问乙出南门直行一百三十五步，甲从乾隅南行六百步，斜望乙与城参相直。问答同前。}
  {假设乙从南门出发向南直走135步，甲从乾隅（西北角）向南行走600步，斜向遥望乙，视线与城池边缘恰好在一条直线上。问题与答案同前。}

\bimethod
  {倍二行差内减甲南行步，复以乘甲南行步为实，从空，倍二行差减甲南行步为廉，一步常法，得城径。}
  {两倍的两行差减去甲向南行走步数，再乘以甲向南行走步数作为被开方数，一次项为空，两倍的两行差减去甲向南行走步数作为二次项系数，三次项系数为1，开方求得城径。}

%% --- Problem 10 ---
\stepcounter{probnum}
\biproblem
  {或问乙出南门直行一百三十五步，甲从乾隅南行六百步望乙。问答同前。}
  {假设乙从南门出发向南直走135步，甲从乾隅（西北角）向南行走600步遥望乙。问题与答案同前。}

\bimethod
  {两行步相乘为实，从空，两行步为廉，一步常法，得城径。}
  {两行步数相乘作为被开方数，一次项为空，两行步数为二次项系数，三次项系数为1，开方求得城径。}

%% --- Problem 11 ---
\stepcounter{probnum}
\biproblem
  {或问乙出南门直行一百三十五步，甲从乾隅南行六百步，斜望乙与城参相直。问答同前。}
  {假设乙从南门出发向南直走135步，甲从乾隅（西北角）向南行走600步，斜向遥望乙，视线与城池边缘恰好在一条直线上。问题与答案同前。}

\bimethod
  {倍二行差内减甲南行步，复以乘甲南行步为实，从空，倍二行差减甲南行步为廉，一步常法，得城径。}
  {两倍的两行差减去甲向南行走步数，再乘以甲向南行走步数作为被开方数，一次项为空，两倍的两行差减去甲向南行走步数作为二次项系数，三次项系数为1，开方求得城径。}

%% --- Problem 12 ---
\stepcounter{probnum}
\biproblem
  {或问乙出南门直行一百三十五步，甲从乾隅南行六百步望乙。问答同前。}
  {假设乙从南门出发向南直走135步，甲从乾隅（西北角）向南行走600步遥望乙。问题与答案同前。}

\bimethod
  {两行步相乘为实，从空，两行步为廉，一步常法，得城径。}
  {两行步数相乘作为被开方数，一次项为空，两行步数为二次项系数，三次项系数为1，开方求得城径。}

%% --- Problem 13 ---
\stepcounter{probnum}
\biproblem
  {或问乙出南门直行一百三十五步，甲从乾隅南行六百步，斜望乙与城参相直。问答同前。}
  {假设乙从南门出发向南直走135步，甲从乾隅（西北角）向南行走600步，斜向遥望乙，视线与城池边缘恰好在一条直线上。问题与答案同前。}

\bimethod
  {倍二行差内减甲南行步，复以乘甲南行步为实，从空，倍二行差减甲南行步为廉，一步常法，得城径。}
  {两倍的两行差减去甲向南行走步数，再乘以甲向南行走步数作为被开方数，一次项为空，两倍的两行差减去甲向南行走步数作为二次项系数，三次项系数为1，开方求得城径。}

%% --- Problem 14 ---
\stepcounter{probnum}
\biproblem
  {或问乙出南门直行一百三十五步，甲从乾隅南行六百步望乙。问答同前。}
  {假设乙从南门出发向南直走135步，甲从乾隅（西北角）向南行走600步遥望乙。问题与答案同前。}

\bimethod
  {两行步相乘为实，从空，两行步为廉，一步常法，得城径。}
  {两行步数相乘作为被开方数，一次项为空，两行步数为二次项系数，三次项系数为1，开方求得城径。}

%% --- Problem 15 ---
\stepcounter{probnum}
\biproblem
  {或问乙出南门直行一百三十五步，甲从乾隅南行六百步，斜望乙与城参相直。问答同前。}
  {假设乙从南门出发向南直走135步，甲从乾隅（西北角）向南行走600步，斜向遥望乙，视线与城池边缘恰好在一条直线上。问题与答案同前。}

\bimethod
  {倍二行差内减甲南行步，复以乘甲南行步为实，从空，倍二行差减甲南行步为廉，一步常法，得城径。}
  {两倍的两行差减去甲向南行走步数，再乘以甲向南行走步数作为被开方数，一次项为空，两倍的两行差减去甲向南行走步数作为二次项系数，三次项系数为1，开方求得城径。}

%% --- Problem 16 ---
\stepcounter{probnum}
\biproblem
  {或问乙出南门直行一百三十五步，甲从乾隅南行六百步望乙。问答同前。}
  {假设乙从南门出发向南直走135步，甲从乾隅（西北角）向南行走600步遥望乙。问题与答案同前。}

\bimethod
  {两行步相乘为实，从空，两行步为廉，一步常法，得城径。}
  {两行步数相乘作为被开方数，一次项为空，两行步数为二次项系数，三次项系数为1，开方求得城径。}

%% --- Problem 17 ---
\stepcounter{probnum}
\biproblem
  {或问乙出南门直行一百三十五步，甲从乾隅南行六百步，斜望乙与城参相直。问答同前。}
  {假设乙从南门出发向南直走135步，甲从乾隅（西北角）向南行走600步，斜向遥望乙，视线与城池边缘恰好在一条直线上。问题与答案同前。}

\bimethod
  {倍二行差内减甲南行步，复以乘甲南行步为实，从空，倍二行差减甲南行步为廉，一步常法，得城径。}
  {两倍的两行差减去甲向南行走步数，再乘以甲向南行走步数作为被开方数，一次项为空，两倍的两行差减去甲向南行走步数作为二次项系数，三次项系数为1，开方求得城径。}

%% --- Problem 18 ---
\stepcounter{probnum}
\biproblem
  {或问乙出南门直行一百三十五步，甲从乾隅南行六百步望乙。问答同前。}
  {假设乙从南门出发向南直走135步，甲从乾隅（西北角）向南行走600步遥望乙。问题与答案同前。}

\bimethod
  {两行步相乘为实，从空，两行步为廉，一步常法，得城径。}
  {两行步数相乘作为被开方数，一次项为空，两行步数为二次项系数，三次项系数为1，开方求得城径。}

\vspace{0.5cm}
\longline
\begin{center}
\textit{测圆海镜卷五终}
\end{center}

\newpage

%% ===================================================================
%% VOLUME 6
%% ===================================================================
\section{测圆海镜卷六 —— 大勾一十八问}
\label{sec:vol6}

\bisectionheader
  {大勾一十八问}
  {十八道关于大勾的数学题}
  {卷六论「大勾」之术。所谓大勾，乃最大勾股形之勾边，即通勾也。以大勾为主，配以诸勾股形之关系，求圆城之径。凡一十八问。}
  {卷六讨论「大勾」的算法。所谓大勾，是指最大勾股三角形的勾边，即通勾。以大勾为主要对象，配合各种勾股三角形的关系，求圆形城池的直径。共18道题。}

\vspace{0.3cm}

%% Reset problem counter for Volume 6
\setcounter{probnum}{0}

%% --- Problem 1 ---
\stepcounter{probnum}
\biproblem
  {或问乙从东门直行一十六步，甲从乾隅东行三百二十步望乙，与城参相直。问答同前。}
  {假设乙从东门出发向东直走16步，甲从乾隅（西北角）向东行走320步遥望乙，视线与城池边缘恰好在一条直线上。问题与答案同前。}

\bian answer
  {城径二百四十步。}
  {城池直径为240步。}

\biworking
  {立天元一为半径，以二之加乙东行，得\raisebox{-0.5ex}{\scriptsize$\begin{array}{c}32\\2\end{array}$}为小差。半乙南行步，得一百八十步，以乘小差，得半径幂，寄左。然后以天元幂与左相消，开平方得一百二十步，即半径。}
  {设天元一为半径，用两倍的天元加上乙向东行走的步数，得到$\begin{pmatrix}32\\2\end{pmatrix}$为小差。取乙向南行走步数的一半，得180步，乘以小差，得到半径的平方，暂存左边。然后将天元的平方与左边相消，开平方得到120步，即半径。}

\bimethod
  {甲东行内减二之乙南行，复以乘甲东行为实。四之东行内减二之乙东行为从，四益隅，得半径。}
  {甲向东行走步数减去两倍乙向南行走步数，再乘以甲向东行走步数作为被开方数。四倍的东行步数减去两倍乙向东行走步数作为一次项系数，四次项采用益积法（增积法），开方求得半径。}

%% --- Problem 2 ---
\stepcounter{probnum}
\biproblem
  {或问乙从东门直行一十六步，甲从乾隅东行三百二十步望乙，与城参相直。问答同前。}
  {假设乙从东门出发向东直走16步，甲从乾隅（西北角）向东行走320步遥望乙，视线与城池边缘恰好在一条直线上。问题与答案同前。}

\bian answer
  {城径二百四十步。}
  {城池直径为240步。}

\biworking
  {立天元一为半径，以二之加乙东行，得\raisebox{-0.5ex}{\scriptsize$\begin{array}{c}16\\2\end{array}$}为小差。半甲东行步，得一百六十步，以乘小差，得\raisebox{-0.5ex}{\scriptsize$\begin{array}{c}2560\\320\end{array}$}为半径幂，寄左。然后以天元幂\raisebox{-0.5ex}{\scriptsize$\begin{array}{c}0\\1\end{array}$}与左相消，开平方得一百二十步，即半径。}
  {设天元一为半径，用两倍的天元加上乙向东行走的步数，得到$\begin{pmatrix}16\\2\end{pmatrix}$为小差。取甲向东行走步数的一半，得160步，乘以小差，得到$\begin{pmatrix}2560\\320\end{pmatrix}$作为半径的平方，暂存左边。然后将天元的平方$\begin{pmatrix}0\\1\end{pmatrix}$与左边相消，开平方得到120步，即半径。}

\bimethod
  {甲东行内减二之乙东行，复以乘甲东行为实。四之甲东行内减二之乙东行为从，四益隅，得半径。}
  {甲向东行走步数减去两倍乙向东行走步数，再乘以甲向东行走步数作为被开方数。四倍的甲向东行走步数减去两倍乙向东行走步数作为一次项系数，四次项采用益积法（增积法），开方求得半径。}

%% --- Problem 3 ---
\stepcounter{probnum}
\biproblem
  {或问乙从东门直行一十六步，甲从乾隅东行三百二十步望乙。问答同前。}
  {假设乙从东门出发向东直走16步，甲从乾隅（西北角）向东行走320步遥望乙。问题与答案同前。}

\bian answer
  {城径二百四十步。}
  {城池直径为240步。}

\biworking
  {立天元一为城径，以二之加乙东行，得\raisebox{-0.5ex}{\scriptsize$\begin{array}{c}16\\2\end{array}$}为小差。置甲东行幂内减小差幂，而半之，得大弦。带小差分母。又置甲东行幂内减小差幂，而半之，得大股。带小差分母。又以小差乘股，併入和，得下式，为小和，以乘大弦。寄左。又以倍天元加斜步，得小和，以乘大弦，得同数，与左相消。开立方得二百四十步，即城径。}
  {设天元一为城池直径，用两倍的天元加上乙向东行走的步数，得到$\begin{pmatrix}16\\2\end{pmatrix}$为小差。将甲向东行走步数的平方减去小差的平方，再取一半，得到大弦。带小差分母。又将甲向东行走步数的平方减去小差的平方，再取一半，得到大股。也带小差分母。又用小差乘以股，合并入和，得到下式，作为小和，乘以大弦。暂存左边。又用两倍的天元加上斜行步数，得到小和，乘以大弦，得到相同的数，与左边相消。开立方得到240步，即城池直径。}

\bimethod
  {甲东行内减二之乙东行，复以乘甲东行为实。四之甲东行内减二之乙东行为从，四益隅，得城径。}
  {甲向东行走步数减去两倍乙向东行走步数，再乘以甲向东行走步数作为被开方数。四倍的甲向东行走步数减去两倍乙向东行走步数作为一次项系数，四次项采用益积法（增积法），开方求得城径。}

\vspace{0.3cm}
\noindent
\begin{minipage}[t]{0.48\textwidth}
\textit{（以下各问草法类推，皆依大勾之术，以小差为主，建立方程。）}
\end{minipage}%
\hfill\vrule\hfill
\begin{minipage}[t]{0.48\textwidth}
\textit{（以下各题演算方法类似，都依照大勾的算法，以小差为主要对象，建立方程。）}
\end{minipage}

\vspace{0.3cm}

%% --- Problem 4 ---
\stepcounter{probnum}
\biproblem
  {或问乙从东门直行一十六步，甲从乾隅东行三百二十步望乙。问答同前。}
  {假设乙从东门出发向东直走16步，甲从乾隅（西北角）向东行走320步遥望乙。问题与答案同前。}

\bimethod
  {甲东行内减二之乙东行，复以乘甲东行为实，从空，四之甲东行内减二之乙东行为廉，四益隅，得城径。}
  {甲向东行走步数减去两倍乙向东行走步数，再乘以甲向东行走步数作为被开方数，一次项为空，四倍的甲向东行走步数减去两倍乙向东行走步数作为二次项系数，四次项采用益积法（增积法），开方求得城径。}

%% --- Problem 5 ---
\stepcounter{probnum}
\biproblem
  {或问乙从东门直行一十六步，甲从乾隅东行三百二十步望乙，与城参相直。问答同前。}
  {假设乙从东门出发向东直走16步，甲从乾隅（西北角）向东行走320步遥望乙，视线与城池边缘恰好在一条直线上。问题与答案同前。}

\bimethod
  {甲东行内减二之乙东行，复以乘甲东行为实，从空，四之甲东行内减二之乙东行为廉，四益隅，得城径。}
  {甲向东行走步数减去两倍乙向东行走步数，再乘以甲向东行走步数作为被开方数，一次项为空，四倍的甲向东行走步数减去两倍乙向东行走步数作为二次项系数，四次项采用益积法（增积法），开方求得城径。}

%% --- Problem 6 ---
\stepcounter{probnum}
\biproblem
  {或问乙从东门直行一十六步，甲从乾隅东行三百二十步望乙。问答同前。}
  {假设乙从东门出发向东直走16步，甲从乾隅（西北角）向东行走320步遥望乙。问题与答案同前。}

\bimethod
  {甲东行内减二之乙东行，复以乘甲东行为实，从空，四之甲东行内减二之乙东行为廉，四益隅，得城径。}
  {甲向东行走步数减去两倍乙向东行走步数，再乘以甲向东行走步数作为被开方数，一次项为空，四倍的甲向东行走步数减去两倍乙向东行走步数作为二次项系数，四次项采用益积法（增积法），开方求得城径。}

%% --- Problem 7 ---
\stepcounter{probnum}
\biproblem
  {或问乙从东门直行一十六步，甲从乾隅东行三百二十步望乙，与城参相直。问答同前。}
  {假设乙从东门出发向东直走16步，甲从乾隅（西北角）向东行走320步遥望乙，视线与城池边缘恰好在一条直线上。问题与答案同前。}

\bimethod
  {甲东行内减二之乙东行，复以乘甲东行为实，从空，四之甲东行内减二之乙东行为廉，四益隅，得城径。}
  {甲向东行走步数减去两倍乙向东行走步数，再乘以甲向东行走步数作为被开方数，一次项为空，四倍的甲向东行走步数减去两倍乙向东行走步数作为二次项系数，四次项采用益积法（增积法），开方求得城径。}

%% --- Problem 8 ---
\stepcounter{probnum}
\biproblem
  {或问乙从东门直行一十六步，甲从乾隅东行三百二十步望乙。问答同前。}
  {假设乙从东门出发向东直走16步，甲从乾隅（西北角）向东行走320步遥望乙。问题与答案同前。}

\bimethod
  {甲东行内减二之乙东行，复以乘甲东行为实，从空，四之甲东行内减二之乙东行为廉，四益隅，得城径。}
  {甲向东行走步数减去两倍乙向东行走步数，再乘以甲向东行走步数作为被开方数，一次项为空，四倍的甲向东行走步数减去两倍乙向东行走步数作为二次项系数，四次项采用益积法（增积法），开方求得城径。}

%% --- Problem 9 ---
\stepcounter{probnum}
\biproblem
  {或问乙从东门直行一十六步，甲从乾隅东行三百二十步望乙，与城参相直。问答同前。}
  {假设乙从东门出发向东直走16步，甲从乾隅（西北角）向东行走320步遥望乙，视线与城池边缘恰好在一条直线上。问题与答案同前。}

\bimethod
  {甲东行内减二之乙东行，复以乘甲东行为实，从空，四之甲东行内减二之乙东行为廉，四益隅，得城径。}
  {甲向东行走步数减去两倍乙向东行走步数，再乘以甲向东行走步数作为被开方数，一次项为空，四倍的甲向东行走步数减去两倍乙向东行走步数作为二次项系数，四次项采用益积法（增积法），开方求得城径。}

%% --- Problem 10 ---
\stepcounter{probnum}
\biproblem
  {或问乙从东门直行一十六步，甲从乾隅东行三百二十步望乙。问答同前。}
  {假设乙从东门出发向东直走16步，甲从乾隅（西北角）向东行走320步遥望乙。问题与答案同前。}

\bimethod
  {甲东行内减二之乙东行，复以乘甲东行为实，从空，四之甲东行内减二之乙东行为廉，四益隅，得城径。}
  {甲向东行走步数减去两倍乙向东行走步数，再乘以甲向东行走步数作为被开方数，一次项为空，四倍的甲向东行走步数减去两倍乙向东行走步数作为二次项系数，四次项采用益积法（增积法），开方求得城径。}

%% --- Problem 11 ---
\stepcounter{probnum}
\biproblem
  {或问乙从东门直行一十六步，甲从乾隅东行三百二十步望乙，与城参相直。问答同前。}
  {假设乙从东门出发向东直走16步，甲从乾隅（西北角）向东行走320步遥望乙，视线与城池边缘恰好在一条直线上。问题与答案同前。}

\bimethod
  {甲东行内减二之乙东行，复以乘甲东行为实，从空，四之甲东行内减二之乙东行为廉，四益隅，得城径。}
  {甲向东行走步数减去两倍乙向东行走步数，再乘以甲向东行走步数作为被开方数，一次项为空，四倍的甲向东行走步数减去两倍乙向东行走步数作为二次项系数，四次项采用益积法（增积法），开方求得城径。}

%% --- Problem 12 ---
\stepcounter{probnum}
\biproblem
  {或问乙从东门直行一十六步，甲从乾隅东行三百二十步望乙。问答同前。}
  {假设乙从东门出发向东直走16步，甲从乾隅（西北角）向东行走320步遥望乙。问题与答案同前。}

\bimethod
  {甲东行内减二之乙东行，复以乘甲东行为实，从空，四之甲东行内减二之乙东行为廉，四益隅，得城径。}
  {甲向东行走步数减去两倍乙向东行走步数，再乘以甲向东行走步数作为被开方数，一次项为空，四倍的甲向东行走步数减去两倍乙向东行走步数作为二次项系数，四次项采用益积法（增积法），开方求得城径。}

%% --- Problem 13 ---
\stepcounter{probnum}
\biproblem
  {或问乙从东门直行一十六步，甲从乾隅东行三百二十步望乙，与城参相直。问答同前。}
  {假设乙从东门出发向东直走16步，甲从乾隅（西北角）向东行走320步遥望乙，视线与城池边缘恰好在一条直线上。问题与答案同前。}

\bimethod
  {甲东行内减二之乙东行，复以乘甲东行为实，从空，四之甲东行内减二之乙东行为廉，四益隅，得城径。}
  {甲向东行走步数减去两倍乙向东行走步数，再乘以甲向东行走步数作为被开方数，一次项为空，四倍的甲向东行走步数减去两倍乙向东行走步数作为二次项系数，四次项采用益积法（增积法），开方求得城径。}

%% --- Problem 14 ---
\stepcounter{probnum}
\biproblem
  {或问乙从东门直行一十六步，甲从乾隅东行三百二十步望乙。问答同前。}
  {假设乙从东门出发向东直走16步，甲从乾隅（西北角）向东行走320步遥望乙。问题与答案同前。}

\bimethod
  {甲东行内减二之乙东行，复以乘甲东行为实，从空，四之甲东行内减二之乙东行为廉，四益隅，得城径。}
  {甲向东行走步数减去两倍乙向东行走步数，再乘以甲向东行走步数作为被开方数，一次项为空，四倍的甲向东行走步数减去两倍乙向东行走步数作为二次项系数，四次项采用益积法（增积法），开方求得城径。}

%% --- Problem 15 ---
\stepcounter{probnum}
\biproblem
  {或问乙从东门直行一十六步，甲从乾隅东行三百二十步望乙，与城参相直。问答同前。}
  {假设乙从东门出发向东直走16步，甲从乾隅（西北角）向东行走320步遥望乙，视线与城池边缘恰好在一条直线上。问题与答案同前。}

\bimethod
  {甲东行内减二之乙东行，复以乘甲东行为实，从空，四之甲东行内减二之乙东行为廉，四益隅，得城径。}
  {甲向东行走步数减去两倍乙向东行走步数，再乘以甲向东行走步数作为被开方数，一次项为空，四倍的甲向东行走步数减去两倍乙向东行走步数作为二次项系数，四次项采用益积法（增积法），开方求得城径。}

%% --- Problem 16 ---
\stepcounter{probnum}
\biproblem
  {或问乙从东门直行一十六步，甲从乾隅东行三百二十步望乙。问答同前。}
  {假设乙从东门出发向东直走16步，甲从乾隅（西北角）向东行走320步遥望乙。问题与答案同前。}

\bimethod
  {甲东行内减二之乙东行，复以乘甲东行为实，从空，四之甲东行内减二之乙东行为廉，四益隅，得城径。}
  {甲向东行走步数减去两倍乙向东行走步数，再乘以甲向东行走步数作为被开方数，一次项为空，四倍的甲向东行走步数减去两倍乙向东行走步数作为二次项系数，四次项采用益积法（增积法），开方求得城径。}

%% --- Problem 17 ---
\stepcounter{probnum}
\biproblem
  {或问乙从东门直行一十六步，甲从乾隅东行三百二十步望乙，与城参相直。问答同前。}
  {假设乙从东门出发向东直走16步，甲从乾隅（西北角）向东行走320步遥望乙，视线与城池边缘恰好在一条直线上。问题与答案同前。}

\bimethod
  {甲东行内减二之乙东行，复以乘甲东行为实，从空，四之甲东行内减二之乙东行为廉，四益隅，得城径。}
  {甲向东行走步数减去两倍乙向东行走步数，再乘以甲向东行走步数作为被开方数，一次项为空，四倍的甲向东行走步数减去两倍乙向东行走步数作为二次项系数，四次项采用益积法（增积法），开方求得城径。}

%% --- Problem 18 ---
\stepcounter{probnum}
\biproblem
  {或问乙从东门直行一十六步，甲从乾隅东行三百二十步望乙，与城参相直。问答同前。}
  {假设乙从东门出发向东直走16步，甲从乾隅（西北角）向东行走320步遥望乙，视线与城池边缘恰好在一条直线上。问题与答案同前。}

\bimethod
  {甲东行内减二之乙东行，复以乘甲东行为实，从空，四之甲东行内减二之乙东行为廉，四益隅，得城径。}
  {甲向东行走步数减去两倍乙向东行走步数，再乘以甲向东行走步数作为被开方数，一次项为空，四倍的甲向东行走步数减去两倍乙向东行走步数作为二次项系数，四次项采用益积法（增积法），开方求得城径。}

\vspace{0.5cm}
\longline
\begin{center}
\textit{测圆海镜卷六终}
\end{center}

\newpage

%% ===================================================================
%% APPENDIX: STRUCTURE OF THE WORK
%% ===================================================================
\section*{附录：全书结构}
\addcontentsline{toc}{section}{附录：全书结构}

\noindent
\begin{minipage}[t]{0.48\textwidth}
\textbf{【原文说明】}

《测圆海镜》全书十二卷，共一百七十问：

\begin{center}
\begin{tabular}{clc}
\toprule
\textbf{卷} & \textbf{篇名} & \textbf{问数} \\
\midrule
卷一 & 圆城图式 & --- \\
卷二 & 正率一十四问 & 14 \\
卷三 & 边股一十七问 & 17 \\
\rowcolor{gray!15}
卷四 & 底勾一十七问 & 17 \\
\rowcolor{gray!15}
卷五 & 大股一十八问 & 18 \\
\rowcolor{gray!15}
卷六 & 大勾一十八问 & 18 \\
卷七 & 明勾一十六问 & 16 \\
卷八 & 明股一十六问 & 16 \\
卷九 & 杂题 & 18 \\
卷十 & 杂题 & 18 \\
卷十一 & 杂题 & 18 \\
卷十二 & 杂题 & 10 \\
\midrule
 & \textbf{合计} & \textbf{170} \\
\bottomrule
\end{tabular}
\end{center}

\vspace{0.5cm}
\textbf{术语解释：}
\begin{itemize}[leftmargin=1em]
  \item \textbf{天元一}：设未知数为一，即现代代数之$x$
  \item \textbf{寄左}：将某式暂存一边，相当于方程之一边
  \item \textbf{相消}：两式相减，相当于移项合并
  \item \textbf{带分母}：某式含有分母未除
  \item \textbf{幂}：平方，即现代记号$x^2$
  \item \textbf{开方}：开平方，即求平方根
  \item \textbf{开立方}：开立方，即求立方根
  \item \textbf{实}：常数项
  \item \textbf{从}：一次项系数
  \item \textbf{廉}：二次项系数
  \item \textbf{隅}：三次项系数（立方方程中）
\end{itemize}
\end{minipage}%
\hfill\vrule\hfill
\begin{minipage}[t]{0.48\textwidth}
\textbf{【今译说明】}

《测圆海镜》全书共12卷，总计170道题：

\begin{center}
\begin{tabular}{clc}
\toprule
\textbf{卷} & \textbf{篇名} & \textbf{题数} \\
\midrule
卷一 & 圆城图式 & --- \\
卷二 & 正率14题 & 14 \\
卷三 & 边股17题 & 17 \\
\rowcolor{gray!15}
卷四 & 底勾17题 & 17 \\
\rowcolor{gray!15}
卷五 & 大股18题 & 18 \\
\rowcolor{gray!15}
卷六 & 大勾18题 & 18 \\
卷七 & 明勾16题 & 16 \\
卷八 & 明股16题 & 16 \\
卷九 & 杂题 & 18 \\
卷十 & 杂题 & 18 \\
卷十一 & 杂题 & 18 \\
卷十二 & 杂题 & 10 \\
\midrule
 & \textbf{合计} & \textbf{170} \\
\bottomrule
\end{tabular}
\end{center}

\vspace{0.5cm}
\textbf{术语对照：}
\begin{itemize}[leftmargin=1em]
  \item \textbf{天元一}：设未知数为一，即现代代数的$x$
  \item \textbf{寄左}：将某算式暂存一边，相当于方程的一边
  \item \textbf{相消}：两式相减，相当于移项合并同类项
  \item \textbf{带分母}：某算式含有分母未除尽
  \item \textbf{幂}：平方，即现代记号$x^2$
  \item \textbf{开方}：开平方，即求平方根
  \item \textbf{开立方}：开立方，即求立方根
  \item \textbf{实}：常数项
  \item \textbf{从}：一次项系数
  \item \textbf{廉}：二次项系数
  \item \textbf{隅}：三次项系数（立方方程中）
\end{itemize}
\end{minipage}

\vspace{1cm}

%% ===================================================================
%% GLOSSARY OF TECHNICAL TERMS
%% ===================================================================
\section*{术语文汇对照}
\addcontentsline{toc}{section}{术语文汇对照}

\noindent
\begin{minipage}[t]{0.48\textwidth}
\textbf{【古代数学术语】}

\begin{tabular}{ll}
\toprule
\textbf{术语} & \textbf{释义} \\
\midrule
天元术 & 设未知数列方程之法 \\
天元一 & 未知数 $x$ \\
勾股 & 直角三角形之两边 \\
勾 & 短直角边 \\
股 & 长直角边 \\
弦 & 斜边 \\
幂 & 平方 $x^2$ \\
实 & 常数项 \\
从 & 一次项系数 \\
廉 & 二次项系数 \\
隅 & 三次项系数 \\
带从开方 & 带参数之方程求解 \\
益积 & 增积法 \\
减从 & 减项法 \\
重差 & 双重差分测量法 \\
三乘方 & 三次方程 \\
四乘方 & 四次方程 \\
寄左 & 暂存一边 \\
相消 & 两式相减 \\
带分母 & 含分母未除 \\
\bottomrule
\end{tabular}
\end{minipage}%
\hfill\vrule\hfill
\begin{minipage}[t]{0.48\textwidth}
\textbf{【现代数学术语】}

\begin{tabular}{ll}
\toprule
\textbf{术语} & \textbf{现代解释} \\
\midrule
天元术 & 代数方法，设未知数列方程 \\
天元一 & 未知数 $x$ \\
勾股 & 直角三角形的两条直角边 \\
勾 & 短直角边（现代记为 $a$） \\
股 & 长直角边（现代记为 $b$） \\
弦 & 斜边（现代记为 $c$） \\
幂 & 平方运算 $x^2$ \\
实 & 方程的常数项 \\
从 & 一次项系数（现代记为 $a_1$） \\
廉 & 二次项系数（现代记为 $a_2$） \\
隅 & 三次项系数（现代记为 $a_3$） \\
带从开方 & 带参数的方程求解（如 $x^2+px=q$） \\
益积 & 增积法，增加被开方数 \\
减从 & 减项法，减少中间项 \\
重差 & 双重差分测量法（三角测量原理） \\
三乘方 & 三次方程（立方方程） \\
四乘方 & 四次方程（双二次方程） \\
寄左 & 暂存方程的一边 \\
相消 & 方程两边相减（移项合并） \\
带分母 & 算式含有未约去的分母 \\
\bottomrule
\end{tabular}
\end{minipage}

\vspace{1cm}
\begin{center}
\rule{0.3\textwidth}{0.4pt}

\textit{--- 全书完 ---}

\vspace{0.5cm}

{\small 文言文原文依《钦定四库全书》本 \\[0.2cm]
白话文翻译与注释为现代译解}
\end{center}

\end{document}
